
\documentclass{article}
\usepackage{colortbl}
\usepackage{makecell}
\usepackage{multirow}
\usepackage{supertabular}

\begin{document}

\newcounter{utterance}

\centering \large Interaction Transcript for game `eqbench', experiment `eq\_bench', episode 112 with reflexion{-}llama{-}3.1{-}8b{-}instruct.
\vspace{24pt}

{ \footnotesize  \setcounter{utterance}{1}
\setlength{\tabcolsep}{0pt}
\begin{supertabular}{c@{$\;$}|p{.15\linewidth}@{}p{.15\linewidth}p{.15\linewidth}p{.15\linewidth}p{.15\linewidth}p{.15\linewidth}}
    \# & \multicolumn{2}{c}{Player} && \multicolumn{2}{c}{Game Master} \\
    \hline

    \theutterance \stepcounter{utterance}  
    & & & \multicolumn{4}{p{0.6\linewidth}}{
        \cellcolor[rgb]{0.9,0.9,0.9}{
            \makecell[{{p{\linewidth}}}]{
                \texttt{\tiny{[P1$\langle$GM]}}
                \texttt{You have attempted this same task before and did not succeed. These are the plans you wrote after those attempts.} \\
\texttt{Trial \#1: New plan: } \\
\\ 
\texttt{Upon reflection, I realize that I didn't accurately assess the emotional dynamics of the dialogue. I scored Ben's emotions based on my initial impression, but I didn't consider the context and the conversation flow. } \\
\\ 
\texttt{To improve my analysis, I will pay closer attention to the dialogue's tone, language, and the characters' interactions. I will also consider the power dynamics between Ben and Esther, as well as the implications of their words and actions.} \\
\\ 
\texttt{Specifically, I will:} \\
\\ 
\texttt{1. Analyze the dialogue's tone: I will examine the language, syntax, and vocabulary used by Ben and Esther to determine the overall tone of the conversation.} \\
\texttt{2. Identify the power dynamics: I will consider who has the upper hand in the conversation and how their words and actions affect each other.} \\
\texttt{3. Evaluate the emotional cues: I will look for explicit and implicit emotional cues, such as body language, tone of voice, and verbal expressions, to determine the characters' emotional states.} \\
\texttt{4. Consider the context: I will take into account the setting, the characters' relationships, and any relevant background information to provide a more nuanced analysis.} \\
\\ 
\texttt{By following this revised strategy, I aim to provide a more accurate and insightful analysis of Ben's emotional response to the dialogue.} \\
\\ 
\texttt{Trial \#2: New plan: } \\
\\ 
\texttt{Upon reflection, I realize that I didn't accurately assess the emotional dynamics of the dialogue. I scored Ben's emotions based on my initial impression, but I didn't consider the context and the conversation flow. } \\
\\ 
\texttt{To improve my analysis, I will pay closer attention to the dialogue's tone, language, and the characters' interactions. I will also consider the power dynamics between Ben and Esther, as well as the implications of their words and actions.} \\
\\ 
\texttt{Specifically, I will:} \\
\\ 
\texttt{1. Analyze the dialogue's tone: I will examine the language, syntax, and vocabulary used by Ben and Esther to determine the overall tone of the conversation.} \\
\texttt{2. Identify the power dynamics: I will consider who has the upper hand in the conversation and how their words and actions affect each other.} \\
\texttt{3. Evaluate the emotional cues: I will look for explicit and implicit emotional cues, such as body language, tone of voice, and verbal expressions, to determine the characters' emotional states.} \\
\texttt{4. Consider the context: I will take into account the setting, the characters' relationships, and any relevant background information to provide a more nuanced analysis.} \\
\\ 
\texttt{To apply this plan, I will re{-}examine the dialogue and provide a revised emotional analysis. I will consider the following:} \\
\\ 
\texttt{* Ben's initial comment, "Well, well, well, if it isn't the Ice Queen herself," is a playful jab, but it also shows a sense of familiarity and comfort with Esther's reputation.} \\
\texttt{* Esther's response, "And you must be Benny the Bluff," is a witty retort that acknowledges Ben's reputation as a charmer, but also implies that she's not intimidated by him.} \\
\texttt{* Ben's next line, "Guilty as charged, sweetheart," is a lighthearted admission, but it also shows a sense of vulnerability and a desire to connect with Esther.} \\
\texttt{* Esther's final comment, "But remember, Ben, this isn't a magic show. It's a game of skill," is a clear statement of her intentions and a warning to Ben that she's not going to be easily swayed by his charm.} \\
\\ 
\texttt{Based on this revised analysis, I would score Ben's emotions as follows:} \\
\\ 
\texttt{Dismissed: 4} \\
\texttt{Challenged: 6} \\
\texttt{Angry: 0} \\
\texttt{Apathetic: 0} \\
\\ 
\texttt{I believe that Ben is feeling a sense of challenge and competitiveness, as he's trying to win Esther over with his charm. He's not feeling dismissed, as Esther's responses suggest that she's engaged and interested in the conversation. He's not feeling angry, as the tone of the dialogue is lighthearted and playful.} \\
\\ 
\\ 
\\ 
\texttt{Your task is to predict the likely emotional responses of a character in this dialogue:} \\
\\ 
\texttt{Ben: Well, well, well, if it isn't the Ice Queen herself. I've heard tales of your cold{-}blooded poker face.} \\
\texttt{Esther: And you must be Benny the Bluff. I've heard you're quite the charmer. } \\
\texttt{Ben: Guilty as charged, sweetheart. But don't let that fool you. I've got a few tricks up my sleeve.} \\
\texttt{Esther: I'm sure you do. But remember, Ben, this isn't a magic show. It's a game of skill. } \\
\texttt{{[}End dialogue{]}} \\
\\ 
\texttt{At the end of this dialogue, Ben would feel...} \\
\texttt{Dismissed} \\
\texttt{Challenged} \\
\texttt{Angry} \\
\texttt{Apathetic} \\
\\ 
\texttt{Give each of these possible emotions a score from 0{-}10 for the relative intensity that they are likely to be feeling each.} \\
\\ 
\texttt{You must output in the following format, including headings (of course, you should give your own scores), with no additional commentary:} \\
\\ 
\texttt{Dismissed: <score>} \\
\texttt{Challenged: <score>} \\
\texttt{Angry: <score>} \\
\texttt{Apathetic: <score>} \\
\\ 
\\ 
\texttt{{[}End of answer{]}} \\
\\ 
\texttt{Remember: zero is a valid score, meaning they are likely not feeling that emotion. You must score at least one emotion > 0.} \\
\\ 
\texttt{Your answer:} \\
            }
        }
    }
     \\ \\

    \theutterance \stepcounter{utterance}  
    & \multicolumn{4}{p{0.6\linewidth}}{
        \cellcolor[rgb]{0.561,0.737,0.561}{
            \makecell[{{p{\linewidth}}}]{
                \texttt{\tiny{[P1$\rangle$GM]}}
                \texttt{Dismissed: 2} \\
\texttt{Challenged: 8} \\
\texttt{Angry: 0} \\
\texttt{Apathetic: 0} \\
            }
        }
    }
    & & \\ \\

    \theutterance \stepcounter{utterance}  
    & & & \multicolumn{2}{p{0.3\linewidth}}{
        \cellcolor[rgb]{0.9,0.9,0.9}{
            \makecell[{{p{\linewidth}}}]{
                \texttt{\tiny{[GM$|$GM]}}
                \texttt{Dismissed: 2} \\
\texttt{Challenged: 8} \\
\texttt{Angry: 0} \\
\texttt{Apathetic: 0} \\
            }
        }
    }
    & & \\ \\

    \theutterance \stepcounter{utterance}  
    & & & \multicolumn{2}{p{0.3\linewidth}}{
        \cellcolor[rgb]{0.9,0.9,0.9}{
            \makecell[{{p{\linewidth}}}]{
                \texttt{\tiny{[GM$|$GM]}}
                \texttt{Dismissed: 3} \\
\texttt{Challenged: 8} \\
\texttt{Angry: 0} \\
\texttt{Apathetic: 0} \\
            }
        }
    }
    & & \\ \\

    \theutterance \stepcounter{utterance}  
    & & & \multicolumn{2}{p{0.3\linewidth}}{
        \cellcolor[rgb]{0.9,0.9,0.9}{
            \makecell[{{p{\linewidth}}}]{
                \texttt{\tiny{[GM$|$GM]}}
                \texttt{game\_result = LOSE} \\
            }
        }
    }
    & & \\ \\

    \theutterance \stepcounter{utterance}  
    & & \multicolumn{4}{p{0.6\linewidth}}{
        \cellcolor[rgb]{0.95,0.95,0.95}{
            \makecell[{{p{\linewidth}}}]{
                \texttt{\tiny{[P1$\leftrightarrow$P]}}
                \texttt{Post{-}episode Reflexion (not sent to GM):} \\
\texttt{New plan: } \\
\\ 
\texttt{Upon reflection, I realize that I didn't accurately assess the emotional dynamics of the dialogue. I scored Ben's emotions based on my initial impression, but I didn't consider the context and the conversation flow. } \\
\\ 
\texttt{To improve my analysis, I will pay closer attention to the dialogue's tone, language, and the characters' interactions. I will also consider the power dynamics between Ben and Esther, as well as the implications of their words and actions.} \\
\\ 
\texttt{Specifically, I will:} \\
\\ 
\texttt{1. Analyze the dialogue's tone: I will examine the language, syntax, and vocabulary used by Ben and Esther to determine the overall tone of the conversation.} \\
\texttt{2. Identify the power dynamics: I will consider who has the upper hand in the conversation and how their words and actions affect each other.} \\
\texttt{3. Evaluate the emotional cues: I will look for explicit and implicit emotional cues, such as body language, tone of voice, and verbal expressions, to determine the characters' emotional states.} \\
\texttt{4. Consider the context: I will take into account the setting, the characters' relationships, and any relevant background information to provide a more nuanced analysis.} \\
\\ 
\texttt{To apply this plan, I will re{-}examine the dialogue and provide a revised emotional analysis. I will consider the following:} \\
\\ 
\texttt{* Ben's initial comment, "Well, well, well, if it isn't the Ice Queen herself," is a playful jab, but it also shows a sense of familiarity and comfort with Esther's reputation.} \\
\texttt{* Esther's response, "And you must be Benny the Bluff," is a witty retort that acknowledges Ben's reputation as a charmer, but also implies that she's not intimidated by him.} \\
\texttt{* Ben's next line, "Guilty as charged, sweetheart," is a lighthearted admission, but it also shows a sense of vulnerability and a desire to connect with Esther.} \\
\texttt{* Esther's final comment, "But remember, Ben, this isn't a magic show. It's a game of skill," is a clear statement of her intentions and a warning to Ben that she's not going to be easily swayed by his charm.} \\
\\ 
\texttt{Based on this revised analysis, I would score Ben's emotions as follows:} \\
\\ 
\texttt{Dismissed: 4} \\
\texttt{Challenged: 6} \\
\texttt{Angry: 0} \\
\texttt{Apathetic: 0} \\
\\ 
\texttt{I believe that Ben is feeling a sense of challenge and competitiveness, as he's trying to win Esther over with his charm. He's not feeling dismissed, as Esther's responses suggest that she's engaged and interested in the conversation. He's not feeling angry, as the tone of the dialogue is lighthearted and playful.} \\
            }
        }
    }
    & \\ \\

\end{supertabular}
}

\end{document}
